% Porteuse : Lynn (brique 3). Sources : docs/protocole.md (qui fait foi), docs/datasets.md,
% data/manifests/, data/splits/. Aucun chiffre n'est écrit ici sans venir de ces fichiers.
\section{Données et protocole expérimental}
\label{sec:donnees}

Cette section décrit les jeux de données, la façon dont ils sont décrits (les manifestes), et
les règles de découpage qui garantissent qu'aucune information du test ne sert à
l'entraînement. Chaque règle porte l'identifiant de la spécification (P pour protocole, R pour
représentation) et correspond à un test automatique du dépôt.

\subsection{Jeux de données et rôle de chacun}
\label{sec:jeux}

Les modèles apprennent sur de la parole anglaise jouée par des acteurs et sur des sons
d'environnement, puis sont mesurés sur trois niveaux de test de plus en plus éloignés de
l'entraînement (P7) : en domaine, hors domaine, et les voix de l'équipe
(tableau~\ref{tab:jeux}). Un jeu de test n'entre jamais en entraînement ni en validation.

\begin{table}[ht]
  \centering
  \small
  \begin{tabularx}{\linewidth}{@{}l>{\raggedright\arraybackslash}X>{\raggedright\arraybackslash}p{2.6cm}r>{\raggedright\arraybackslash}p{3.3cm}@{}}
    \toprule
    Jeu & Contenu & Tâche & Fichiers & Rôle \\
    \midrule
    CREMA-D    & 91 acteurs anglophones (48 H, 43 F), 12 phrases & Émotions & 7\,442 & entraînement, validation, test \\
    RAVDESS    & 24 acteurs anglophones (12 H, 12 F), 2 phrases  & Émotions & 1\,440 & entraînement, validation, test \\
    ESC-50     & 50 classes de sons, extraits de 5\,s             & Événements & 2\,000 & validation croisée, 5 folds \\
    CaFE       & 12 acteurs francophones (6 H, 6 F)             & Émotions & 936    & test hors domaine \\
    OpenSLR~57 & français à l'accent africain, 19,9\,h          & Transcription, locuteur & 16\,656 & test seulement \\
    Équipe (prévu) & 4 personnes, téléphone et ordinateur : 20 phrases lues, 2 phrases jouées avec les 6 émotions, 1 ou 2 sons & Toutes & --- & test et démonstration \\
    \bottomrule
  \end{tabularx}
  \caption{Jeux de données et rôle de chacun (P7, P10). La ligne « Équipe » décrit ce que
    prévoit le protocole d'enregistrement (P13) : ces voix ne sont pas encore enregistrées. ACE
    (voix africaines, émotions) était prévu en test hors domaine ; l'accès n'a pas été obtenu.}
  \label{tab:jeux}
\end{table}

\aRediger{Lynn}{Ligne « Équipe » : remplacer la prévision par le nombre réel d'extraits une
  fois les enregistrements faits (T6), ou retirer la ligne. Statut final d'ACE.}

\subsection{Manifestes}
\label{sec:manifestes}

Chaque jeu est décrit par un \emph{manifeste} : un fichier CSV, versionné dans
\texttt{data/manifests/}, qui contient une ligne par fichier audio et aucun son. Ses colonnes
sont le jeu, le chemin relatif du fichier, l'étiquette d'origine, le locuteur, son genre, la
durée, le fold (ESC-50), le découpage officiel et le sous-corpus (OpenSLR~57), l'intensité
jouée et la transcription brute.

Les manifestes sont construits sur Kaggle, à partir des archives officielles, par les
fonctions de \texttt{echolens.core.data} : les informations sont lues dans les \emph{noms} de
fichiers (par exemple \texttt{01-C-1-1.wav} pour CaFE : acteur 01, colère, intensité faible,
phrase 1), jamais dans les noms de dossiers, dont les accents peuvent être mal encodés.
Chaque manifeste est ensuite contrôlé : un seul jeu, des chemins uniques, et un même
enregistrement d'origine jamais réparti sur deux folds (P3). Les effectifs obtenus sont comparés
à ceux attendus, et la construction s'arrête au moindre écart.

Le manifeste garde l'étiquette d'origine (P9) ; la correspondance vers les classes communes se
fait ensuite dans le code. On peut ainsi la modifier sans reconstruire les manifestes.

\subsection{Correspondance des émotions}
\label{sec:correspondance}

On garde six classes présentes dans tous les jeux : neutre, joie, tristesse, colère, peur et
dégoût (P2). Le \emph{calme} de RAVDESS est exclu plutôt que fusionné avec le neutre : c'est une
consigne de jeu à part, et la fusion triplerait les neutres de RAVDESS (de 96 à 288). La
\emph{surprise} est exclue car CREMA-D, le jeu principal, n'en contient pas.

\begin{table}[ht]
  \centering
  \small
  \begin{tabular}{@{}lrrrrr@{}}
    \toprule
    Jeu & Fichiers & Exclus & Gardés & Neutre & Chacune des 5 autres \\
    \midrule
    CREMA-D & 7\,442 & 0   & 7\,442 & 1\,087 & 1\,271 \\
    RAVDESS & 1\,440 & 384 & 1\,056 & 96     & 192 \\
    CaFE    & 936    & 144 & 792    & 72     & 144 \\
    \bottomrule
  \end{tabular}
  \caption{Effectifs après correspondance vers les six classes communes, vérifiés sur les
    manifestes réels.}
  \label{tab:emotions}
\end{table}

Le neutre est un peu moins fréquent que les autres classes : après le découpage par acteur
(section~\ref{sec:p1}), l'entraînement compte 815 neutres contre 1\,007 pour chacune des cinq
autres émotions, la validation et le test 184 contre 228. La métrique principale est le F1
macro (P4), où chaque classe pèse autant quelle que soit sa fréquence ; l'accuracy est aussi
rapportée, pour rester comparable aux articles.

\subsection{Découpage des émotions par acteur}
\label{sec:p1}

Les \emph{acteurs}, et non les fichiers, sont répartis entre entraînement, validation et test
(P1). Si un même acteur apparaissait dans l'entraînement et le test, le modèle pourrait
reconnaître sa voix plutôt que l'émotion, et le score serait gonflé.

Pour un jeu de $A$ acteurs, la validation et le test reçoivent chacun
\begin{equation}
  n = \operatorname{round}\!\left(\frac{0{,}15 \cdot A}{2}\right)
\end{equation}
acteurs de \emph{chaque} genre ; tous les autres vont à l'entraînement. Le tirage est fait une
seule fois, avec une graine fixe (42), et le résultat est versionné dans
\texttt{data/splits/emotion\_split.csv} : tous les modèles utilisent exactement les mêmes
acteurs (tableau~\ref{tab:p1}).

\begin{table}[ht]
  \centering
  \small
  \begin{tabular}{@{}lllll@{}}
    \toprule
    Jeu & Acteurs & Entraînement & Validation & Test \\
    \midrule
    CREMA-D & 91 (48 H, 43 F) & 63 (34 H, 29 F) & 14 (7 H, 7 F) & 14 (7 H, 7 F) \\
    RAVDESS & 24 (12 H, 12 F) & 16 (8 H, 8 F)   & 4 (2 H, 2 F)  & 4 (2 H, 2 F) \\
    \bottomrule
  \end{tabular}
  \caption{Découpage par acteur (P1). Le test compte 1\,324 extraits et 18 acteurs, autant de
    femmes que d'hommes, ce qui permet d'analyser les biais par genre (E4).}
  \label{tab:p1}
\end{table}

\subsection{ESC-50 : les folds officiels}
\label{sec:p3}

Les 2\,000 extraits d'ESC-50 sont des fragments d'enregistrements plus longs, rangés par
l'auteur du jeu dans cinq folds de 400 extraits. On garde ces folds tels quels (P3) : les
remélanger placerait deux fragments du même enregistrement en entraînement et en test. Ces
folds sont aussi ceux de la littérature, ce qui rend nos scores comparables. Quatre
enregistrements (9 extraits) sont, dans le jeu officiel, répartis sur deux folds ; ils sont
listés dans le code, et toute autre source à cheval sur deux folds est refusée.

Au tour $k$ ($k = 1, \dots, 5$), le fold $k$ sert de test. Les réglages (nombre d'époques,
hyperparamètres) sont choisis en entraînant sur trois folds et en validant sur le fold
\begin{equation}
  v = (k \bmod 5) + 1 ,
\end{equation}
puis le modèle est réentraîné sur les quatre folds hors test et évalué une seule fois sur le
fold $k$. On rapporte la moyenne et l'écart-type sur les cinq tours.

\subsection{Standardisation sur l'entraînement seul}
\label{sec:r4}

Chaque bande Mel est centrée et réduite avec des statistiques calculées sur l'entraînement
seulement (R4). Pour la bande $b$,
\begin{equation}
  \mu_b = \frac{1}{N T} \sum_{n=1}^{N} \sum_{t=1}^{T} X_{n,b,t} , \qquad
  \sigma_b = \sqrt{\frac{1}{N T} \sum_{n=1}^{N} \sum_{t=1}^{T} \left(X_{n,b,t} - \mu_b\right)^2} , \qquad
  \tilde{X}_{n,b,t} = \frac{X_{n,b,t} - \mu_b}{\sigma_b} ,
\end{equation}
où $X_{n,b,t}$ est le log-Mel (en dB) de l'extrait $n$ de l'entraînement, à la bande $b$ et à
la trame $t$, $N$ le nombre d'extraits d'entraînement et $T$ le nombre de trames. Les mêmes
$\mu_b$ et $\sigma_b$ sont ensuite appliqués à la validation, au test et à l'application.
Calculer ces statistiques sur le test serait une fuite : le modèle bénéficierait d'une
information sur les données qui servent à l'évaluer. De même, l'augmentation de données (P6)
n'est appliquée qu'après le découpage, et à l'entraînement seul.

\subsection{Identification du locuteur}
\label{sec:p11}

Pour reconnaître une voix, il faut l'avoir entendue : on ne sépare donc pas les locuteurs,
mais leurs \emph{phrases} (P11). Les découpages officiels d'OpenSLR~57 ne conviennent pas, et
ne séparent d'ailleurs pas tous les locuteurs : 14 dossiers du sous-corpus \texttt{ca16} sont
à la fois dans \texttt{train} et dans \texttt{devtest}, avec le même préfixe de session. Un
locuteur est donc identifié par son sous-corpus et son dossier ; 225 locuteurs sur 233 ont au
moins dix phrases et sont gardés. Pour chacun, cinq phrases tirées au hasard (graine 42)
servent à enregistrer sa voix, les autres servent au test.

Les locuteurs sont ensuite répartis en quatre groupes, un sur quatre dans chaque sous-corpus
(tableau~\ref{tab:p11}). Le seuil « locuteur inconnu » est choisi sur le groupe de calibration
et ses imposteurs, puis mesuré sur le groupe d'évaluation et d'autres imposteurs : aucune voix
ne sert à la fois à choisir le seuil et à le mesurer. Comme dans l'application, la décision
compare au seuil le \emph{meilleur} score d'une phrase parmi toutes les voix enregistrées.

\begin{table}[ht]
  \centering
  \small
  \begin{tabular}{@{}llrr@{}}
    \toprule
    Groupe & Rôle & Locuteurs & Phrases de test \\
    \midrule
    Calibration               & voix enregistrées, choix du seuil & 57 & 3\,570 \\
    Imposteurs de calibration & voix inconnues, choix du seuil    & 56 & 3\,906 \\
    Évaluation                & voix enregistrées, mesure         & 57 & 4\,246 \\
    Imposteurs d'évaluation   & voix inconnues, mesure            & 55 & 3\,770 \\
    \bottomrule
  \end{tabular}
  \caption{Groupes de locuteurs d'OpenSLR~57 (P11), versionnés dans
    \texttt{data/splits/openslr57\_speakers.csv}.}
  \label{tab:p11}
\end{table}

Rien ne garantit que deux dossiers différents correspondent à deux personnes différentes :
OpenSLR~57 ne fournit pas d'identité. Si une même personne apparaissait sous deux dossiers, le
score serait sous-estimé.

\subsection{Transcription}
\label{sec:p12}

Le taux d'erreur sur les mots (WER) est mesuré sur le test officiel d'OpenSLR~57 (515 phrases)
et sur le sous-corpus de Yaoundé (8\,397 phrases), dont toutes les transcriptions sont
disponibles (P10). Référence et transcription sont normalisées de la même façon avant la
mesure (P12) : minuscules, ponctuation retirée sauf l'apostrophe et le trait d'union dans un
mot, apostrophes unifiées, nombres écrits en lettres, espaces multiples réduits.

\aRediger{Lynn}{Après C17 : nombre de phrases réellement évaluées et éventuelles exclusions.}

\subsection{Tests de fuite}
\label{sec:tests-fuite}

Chaque règle de cette section est vérifiée par un test automatique, lancé à chaque
modification du code : aucun acteur dans deux découpages d'émotions ; chaque fichier gardé dans
un seul découpage ; toutes les étiquettes dans les six classes après correspondance ; effectifs
par découpage et par genre conformes au tableau~\ref{tab:p1} ; découpage identique quand on le
régénère ; 400 extraits par fold d'ESC-50 ; aucune phrase d'un locuteur à la fois en
enregistrement et en test ; seuil du locuteur inchangé quand on modifie les imposteurs
d'évaluation.

\subsection{Limites des données}
\label{sec:limites-donnees}

\begin{itemize}
  \item \textbf{Phrases communes.} CREMA-D utilise les mêmes 12 phrases pour tous les acteurs,
        RAVDESS les mêmes 2 : le texte est commun à l'entraînement et au test. CaFE et les voix
        de l'équipe apportent des phrases nouvelles.
  \item \textbf{Intensités déséquilibrées dans CREMA-D.} 6\,077 extraits en intensité non
        précisée contre 455 pour chacune des trois autres : une analyse par intensité ne porte
        que sur la seule phrase jouée aux trois niveaux.
  \item \textbf{Repères humains.} Sur l'audio seul de CREMA-D, des auditeurs humains reconnaissent
        l'émotion voulue dans 40,9\,\% des cas ; sur ESC-50, 81,3\,\%. Le hasard vaut
        16,7\,\% pour six émotions et 2\,\% pour cinquante sons.
  \item \textbf{16\,kHz.} Tout ce qui dépasse 8\,kHz est perdu, ce qui pénalise certains sons.
  \item \textbf{Genre dans OpenSLR~57.} Il n'est connu que pour les 12 locuteurs du test
        officiel ; l'analyse des biais par genre y est donc très limitée.
  \item \textbf{Transcriptions incomplètes.} Une partie du sous-corpus \texttt{ca16} n'a pas de
        transcription dans l'archive (devtest 481 sur 761, entraînement 3\,063 sur 5\,997) ; ces
        phrases ne servent pas au WER.
  \item \textbf{Voix de l'équipe.} Quatre locuteurs non acteurs, enregistrés chacun de son
        côté : les résultats sur ces voix sont illustratifs, pas statistiquement concluants.
\end{itemize}
